pdfoutput=1
\pdfoutput=1
% ===========================================================================
%  Audit Instantons: Non-Perturbative Topological Defects in Expert Auditing
%  Instanton
% ===========================================================================
%  A rigorous mathematical formulation inspired by worldsheet instantons 
%  in string theory, applied to the SCX (Situs Consensus eXpert) framework.
%  InstantonSCX
% ===========================================================================

\documentclass[11pt,a4paper]{article}

% ===========================================================================
%  Packages
% ===========================================================================
\usepackage[T1]{fontenc}
\usepackage{amsmath,amssymb,amsthm}
\usepackage{mathtools}
\usepackage{mathrsfs}
\usepackage{bbm}
\usepackage{bm}
\usepackage{graphicx}
\usepackage{xcolor}
\usepackage{hyperref}
\usepackage{enumitem}
\usepackage{booktabs}
\usepackage{tikz}
\usepackage{tikz-cd}
\usepackage{float}
\usepackage{caption}
\usepackage{subcaption}
\usepackage{geometry}
\usepackage[colorinlistoftodos]{todonotes}

% Page geometry
\geometry{margin=1.2in,top=1in,bottom=1in}

% ===========================================================================
%  Theorem environments
% ===========================================================================
\theoremstyle{plain}
\newtheorem{theorem}{Theorem}[section]
\newtheorem{lemma}[theorem]{Lemma}
\newtheorem{proposition}[theorem]{Proposition}
\newtheorem{corollary}[theorem]{Corollary}
\newtheorem{conjecture}[theorem]{Conjecture}

\theoremstyle{definition}
\newtheorem{definition}[theorem]{Definition}
\newtheorem{example}[theorem]{Example}
\newtheorem{remark}[theorem]{Remark}
\newtheorem{protocol}[theorem]{Protocol}

% Custom theorem style for bilingual statements
\newenvironment{bilingual}[2]{%
  \medskip\noindent\textbf{#1}\par\noindent\textit{#2}\par\smallskip
}{}

% ===========================================================================
%  Math operators and notation
% ===========================================================================
\DeclareMathOperator{\im}{im}
\DeclareMathOperator{\coker}{coker}
\DeclareMathOperator{\Hom}{Hom}
\DeclareMathOperator{\Ext}{Ext}
\DeclareMathOperator{\Tor}{Tor}
\DeclareMathOperator{\rank}{rank}
\DeclareMathOperator{\vol}{vol}
\DeclareMathOperator{\diam}{diam}
\DeclareMathOperator{\supp}{supp}
\DeclareMathOperator{\Pers}{Pers}
\DeclareMathOperator{\Birth}{Birth}
\DeclareMathOperator{\Death}{Death}
\DeclareMathOperator{\PH}{PH}
\DeclareMathOperator{\Yajie}{Yajie}
\DeclareMathOperator{\Cercis}{Cercis}
\DeclareMathOperator{\Spec}{Spec}

\newcommand{\R}{\mathbb{R}}
\newcommand{\C}{\mathbb{C}}
\newcommand{\Z}{\mathbb{Z}}
\newcommand{\N}{\mathbb{N}}
\newcommand{\Q}{\mathbb{Q}}
\newcommand{\T}{\mathbb{T}}
\newcommand{\cM}{\mathcal{M}}
\newcommand{\cS}{\mathcal{S}}
\newcommand{\cF}{\mathcal{F}}
\newcommand{\cA}{\mathcal{A}}
\newcommand{\cI}{\mathcal{I}}
\newcommand{\cD}{\mathcal{D}}
\newcommand{\cP}{\mathcal{P}}
\newcommand{\cH}{\mathcal{H}}
\newcommand{\cC}{\mathcal{C}}
\newcommand{\cL}{\mathcal{L}}
\newcommand{\cO}{\mathcal{O}}
\newcommand{\cX}{\mathcal{X}}
\newcommand{\cG}{\mathcal{G}}
\newcommand{\cB}{\mathcal{B}}
\newcommand{\cV}{\mathcal{V}}
\newcommand{\cK}{\mathcal{K}}

% For homology/cohomology groups
\newcommand{\Ho}{H}
\newcommand{\dpartial}{\partial}
\newcommand{\ddelta}{\delta}

% For discrete Hodge
\newcommand{\dH}{d}        % discrete exterior derivative
\newcommand{\dHc}{\delta}  % discrete codifferential

% ===========================================================================
%  Title and Author
% ===========================================================================
\title{
  {\large\textbf{Audit Instantons:}}\\[4pt]
  {\large\textbf{Non-Perturbative Topological Defects in Expert Auditing}}\\[8pt]
  {\normalsize\textit{Instanton}}\\[4pt]
  {\footnotesize A Mathematical Formalization Inspired by String Theory Worldsheet Instantons}\\[2pt]
  {\footnotesize\textit{Instanton}}
}

\author{SCX}
\date{}

% ===========================================================================
\begin{document}
\maketitle

% ===========================================================================
%  Abstract
% ===========================================================================
\begin{abstract}
\noindent
We introduce the concept of \textbf{audit instantons} — non-perturbative, 
topologically protected failures of expert consensus in the SCX (Situs 
Consensus eXpert) auditing framework. Drawing structural inspiration from 
worldsheet instantons in string theory, we demonstrate that pointwise 
Yajie consensus auditing is analogous to perturbative expansion: it 
systematically misses globally correlated expert failures concentrated 
in low-data-density regions of the Situs manifold.

We provide a rigorous mathematical formulation: the Situs manifold 
$\cM$ equipped with a data density function $\rho(x)$; low-density 
sublevel sets $\cM_{\varepsilon} = \{x \in \cM : \rho(x) < \varepsilon\}$; 
and an \emph{audit instanton} defined as a non-contractible 1-cycle 
$\gamma \subset \cM$ on which all experts deviate consistently in the 
same direction — locally exact ($d_1 A = 0$ on $\gamma$) but globally 
non-trivial ($\gamma \notin \im(d_0)$). Detection proceeds via persistent 
homology of the density filtration, where long-lived 1-cycles in the 
persistence diagram correspond to audit instanton candidates. 

We prove the main structural theorem: every non-trivial 1-homology class 
of the Situs manifold contains at least one audit instanton candidate, 
provided the expert output field satisfies mild regularity conditions. 
We further show that audit instantons are invisible to the standard Yajie 
consensus score — precisely the ``non-perturbative'' character that makes 
them dangerous. A concrete detection protocol is given.

\bigskip
\noindent
\textit{we introduce\textbf{Instanton}audit instantons——SCXSitusInstantonProofYajieMissingSitus}

\textit{ $\rho(x)$  Situs  $\cM$ $\cM_{\varepsilon} = \{x \in \cM : \rho(x) < \varepsilon\}$\textbf{Instanton}——define $\cM$  1- $\gamma$all——$d_1 A = 0$  $\gamma$ $\gamma \notin \im(d_0)$where 1-Instanton}

\textit{ProoftheoremsatisfyingSituseach1-InstantonProofInstantonYajieScore——``''}

\bigskip
\noindent
\textbf{Keywords:} audit instantons, persistent homology, discrete Hodge theory, 
topological data analysis, expert auditing, non-perturbative effects, 
SCX framework

\noindent
\textbf{Keywords} InstantonHodgeSCX
\end{abstract}

% ===========================================================================
\tableofcontents

% ===========================================================================
%  SECTION 1: INTRODUCTION
% ===========================================================================
\section{Introduction}
\section*{}

% ---------------------------------------------------------------------------
\subsection{Motivation from Physics}
\subsection*{}

In quantum field theory and string theory, \emph{instantons} are non-perturbative 
field configurations: classical solutions to the Euclidean equations of motion 
with finite action that contribute terms of order $\sim \exp(-1/g^2)$ to the 
path integral, where $g$ is the coupling constant \cite{coleman1977, polyakov1977}. 
Crucially, such contributions are identically zero in any finite-order Taylor 
expansion around $g=0$ — they are \emph{invisible to perturbation theory}.

\bigskip
\noindent
\textit{\textbf{Instanton}instantonsScale $\sim \exp(-1/g^2)$ where $g$  $g=0$ ——\textbf{}}

In string theory specifically, \textbf{worldsheet instantons} arise from string 
worldsheets wrapping holomorphic curves in the compactification manifold. Their 
contribution is $\sim \exp(-A/\alpha')$, where $A$ is the area of the wrapped 
curve and $\alpha'$ is the string tension. These effects are:
\begin{enumerate}[label=(\roman*)]
  \item \textbf{Non-perturbative}: invisible to the genus expansion (Feynman diagrams);
  \item \textbf{Topological}: classified by the homotopy/homology of the target space;
  \item \textbf{Exponentially suppressed} in the small-$g$ regime, yet dominant 
        in certain corners of moduli space.
\end{enumerate}

\bigskip
\noindent
\textit{in particular\textbf{Instanton} $\sim \exp(-A/\alpha')$where $A$ $\alpha'$ (i) \textbf{}Feynman(ii) \textbf{}/(iii)  $g$ \textbf{}}

% ---------------------------------------------------------------------------
\subsection{The Auditing Analogy}
\subsection*{}

The SCX (Situs Consensus eXpert) framework \cite{scx_framework} performs 
expert auditing on a data manifold $\cM$ using discrete Hodge theory and 
gauge-theoretic consensus measures. The standard Yajie audit operates 
\emph{pointwise}: at each data point $x \in \cM$, it compares expert 
outputs and assigns a consensus score based on local agreement. This is 
structurally analogous to \emph{perturbative expansion} in quantum field 
theory — each data point is treated independently, and the audit result is 
the sum of local contributions.

\bigskip
\noindent
\textit{SCXSitusHodge $\cM$ Yajie\textbf{} $\cM$ each $x$ConsistencyScore\textbf{}——each}

The central insight of this paper is that \textbf{pointwise audit is perturbative 
audit}, and it suffers from the same fundamental blindness: it cannot detect 
\emph{globally correlated, topologically protected} failure modes that we term 
\textbf{audit instantons}. These are configurations where:

\begin{enumerate}[label=(\arabic*)]
  \item The data density $\rho(x)$ is low (the ``weak-coupling'' regime of audit);
  \item All experts deviate in the \emph{same} direction — producing \emph{spurious 
        consensus} that looks like high agreement to pointwise Yajie;
  \item The deviation field is locally consistent (closed form) but globally 
        non-trivial (not exact) — a topological obstruction invisible to 
        local analysis.
\end{enumerate}

\bigskip
\noindent
\textit{Core\textbf{i.e.}\textbf{Instanton}\textbf{}satisfying(1)  $\rho(x)$ ``''(2) all\textbf{}——YajieConsistency\textbf{}(3) ——}

% ---------------------------------------------------------------------------
\subsection{Our Contributions}
\subsection*{}

This paper provides the first rigorous mathematical formalization of audit 
instantons. Specifically, we:

\begin{enumerate}[label=(\textbf{C\arabic*})]
  \item \textbf{Formal Definition} (Section 3): Define the Situs manifold 
        $(\cM, \rho, \mathbf{A})$ with data density $\rho$ and audit vector 
        field $\mathbf{A}$, and give a cohomological definition of audit 
        instantons as non-contractible 1-cycles with locally exact but 
        globally non-trivial deviation fields.
  
  \item \textbf{Detection Theory} (Section 4): Show that persistent homology 
        of the density sublevel filtration provides a systematic detection 
        mechanism: 1-cycles with large persistence (long lifetime in the 
        barcode) are audit instanton candidates.
  
  \item \textbf{Structural Theorem} (Theorem \ref{thm:main}): Prove that 
        every non-trivial 1-homology class in $\cM$ contains at least one 
        audit instanton candidate, under mild conditions on the expert 
        output field.
  
  \item \textbf{Invisibility to Yajie} (Section 5): Prove that the standard 
        Yajie consensus score evaluated on an audit instanton cycle 
        systematically overestimates audit quality — the very definition 
        of a non-perturbative effect.
  
  \item \textbf{Detection Protocol} (Section 6): Provide a concrete, 
        algorithmically implementable protocol for detecting audit instantons 
        in real SCX deployments.
\end{enumerate}

\bigskip
\noindent
\textit{This paper provides the firstInstanton\textbf{(C1)} define3define $\rho$  $\mathbf{A}$ Situs $(\cM, \rho, \mathbf{A})$Instantondefine\textbf{(C2)} 4Proof——1-Instanton\textbf{(C3)} theoremtheorem\ref{thm:main}Proof$\cM$ each1-Instanton\textbf{(C4)} Yajie5ProofYajieScoreInstanton——define\textbf{(C5)} 6Instanton}

% ===========================================================================
%  SECTION 2: MATHEMATICAL PRELIMINARIES
% ===========================================================================
\section{Mathematical Preliminaries}
\section*{}

\subsection{The Situs Manifold}
\subsection*{Situs}

\begin{definition}[Situs Manifold / Situs]
\label{def:situs}
Let $\cM$ be a compact, connected, smooth Riemannian $n$-manifold with 
boundary $\partial \cM$ (possibly empty). The \textbf{Situs manifold} is 
the triple $(\cM, \rho, \cE)$, where:
\begin{enumerate}[label=(\roman*)]
  \item $\rho: \cM \to \R_{\geq 0}$ is a smooth (or at least $C^2$) 
        \textbf{data density function}, normalized so that 
        $\int_{\cM} \rho(x) \, d\vol_{\cM}(x) = 1$. 
        In practice, $\rho$ is estimated from the empirical distribution 
        of training or audit data points via kernel density estimation.
  
  \item $\cE = \{E_1, E_2, \ldots, E_M\}$ is a finite set of $M \geq 2$ 
        \textbf{experts}, where each $E_i: \cM \to \R^d$ is a vector-valued 
        function representing the expert's output (prediction, classification, 
        or scalar field) at each point $x \in \cM$.
\end{enumerate}
The \textbf{dimension} $n$ of $\cM$ is typically the intrinsic dimension of 
the data representation space (e.g., after dimensionality reduction).
\end{definition}

\bigskip
\noindent
\textit{\textbf{define \ref{def:situs}Situs} $\cM$  $n$  Riemann  $\partial \cM$\textbf{Situs} $(\cM, \rho, \cE)$where(i) $\rho: \cM \to \R_{\geq 0}$  $C^2$\textbf{}such that $\int_{\cM} \rho(x) \, d\vol_{\cM}(x) = 1$ $\rho$ (ii) $\cE = \{E_1, \ldots, E_M\}$  $M \geq 2$ \textbf{}each $E_i: \cM \to \R^d$ each $x \in \cM$ }

\begin{remark}[Discrete vs. Continuous / ]
In SCX practice, $\cM$ is discretized as a simplicial complex $K$ (or 
graph/cell complex) with vertices corresponding to data points. All 
definitions in this paper extend naturally to the discrete setting via 
discrete Hodge theory \cite{scx_hodge}. We work in the continuous category 
for clarity, but all constructions have well-defined discrete analogues.
\end{remark}

\bigskip
\noindent
\textit{\textbf{remark}SCX$\cM$  $K$/alldefineHodgealldefine}

% ---------------------------------------------------------------------------
\subsection{Data Density and Sublevel Sets}
\subsection*{}

\begin{definition}[Density Sublevel Sets / ]
\label{def:sublevel}
For any threshold $\varepsilon \geq 0$, the \textbf{$\varepsilon$-sublevel set} 
of the data density is:
\begin{equation}
  \cM_{\varepsilon} = \{x \in \cM : \rho(x) \leq \varepsilon\}.
  \label{eq:sublevel}
\end{equation}
By convention, $\cM_0 = \{x : \rho(x) = 0\}$ and $\cM_{\infty} = \cM$.

The family $\{\cM_{\varepsilon}\}_{\varepsilon \geq 0}$ forms a 
\textbf{filtration} of $\cM$: 
$\cM_0 \subseteq \cM_{\varepsilon_1} \subseteq \cM_{\varepsilon_2} \subseteq \cM$ 
for $0 \leq \varepsilon_1 \leq \varepsilon_2$.

The \textbf{critical density} $\rho_{\text{crit}}$ is a distinguished 
threshold below which expert behavior is deemed unreliable due to data 
sparsity. Formally, $\rho_{\text{crit}}$ may be defined as:
\begin{equation}
  \rho_{\text{crit}} = \inf\{\varepsilon > 0 : \text{expert predictions on } 
    \cM_{\varepsilon} \text{ are statistically indistinguishable from 
    their behavior on } \cM \setminus \cM_{\varepsilon}\}.
  \label{eq:rhocrit}
\end{equation}
\end{definition}

\bigskip
\noindent
\textit{\textbf{define \ref{def:sublevel}}any $\varepsilon \geq 0$\textbf{$\varepsilon$-} $\cM_{\varepsilon} = \{x \in \cM : \rho(x) \leq \varepsilon\}$ $\{\cM_{\varepsilon}\}_{\varepsilon \geq 0}$  $\cM$ \textbf{}\textbf{} $\rho_{\text{crit}}$ define $\cM_{\varepsilon}$  $\cM \setminus \cM_{\varepsilon}$  $\varepsilon$}

% ---------------------------------------------------------------------------
\subsection{Discrete Hodge Theory (Recap)}
\subsection*{Hodge}

We briefly recall the essential elements of discrete Hodge theory on $\cM$, 
as used in SCX \cite{scx_hodge}. Let $\Omega^k(\cM)$ denote the space of 
smooth differential $k$-forms on $\cM$.

\begin{definition}[Cochain Complex / ]
The \textbf{de Rham complex} on $\cM$ is:
\begin{equation}
  0 \longrightarrow \Omega^0(\cM) \xrightarrow{d_0} \Omega^1(\cM) 
    \xrightarrow{d_1} \Omega^2(\cM) \xrightarrow{d_2} \cdots 
    \xrightarrow{d_{n-1}} \Omega^n(\cM) \longrightarrow 0,
  \label{eq:derham}
\end{equation}
where each $d_k: \Omega^k(\cM) \to \Omega^{k+1}(\cM)$ is the exterior 
derivative, satisfying $d_{k+1} \circ d_k = 0$.

The \textbf{$k$-th de Rham cohomology} is 
$H^k_{\text{dR}}(\cM) = \ker(d_k) / \im(d_{k-1})$.

The \textbf{Hodge Laplacian} on $k$-forms is:
\begin{equation}
  \Delta_k = d_{k-1} \delta_k + \delta_{k+1} d_k,
  \label{eq:hodge_laplacian}
\end{equation}
where $\delta_k = (-1)^{n(k-1)+1} \star d_{n-k} \star$ is the codifferential 
and $\star$ is the Hodge star operator induced by the Riemannian metric.
\end{definition}

\bigskip
\noindent
\textit{\textbf{define}$\cM$ \textbf{de Rham}(\ref{eq:derham})where $d_k$ satisfying $d_{k+1} \circ d_k = 0$ $k$ de Rham $H^k_{\text{dR}}(\cM) = \ker(d_k) / \im(d_{k-1})$\textbf{Hodge Laplacian} (\ref{eq:hodge_laplacian})where $\delta_k$ $\star$ Hodge}

\begin{definition}[Expert Output 1-Form / 1-]
Given $M$ experts and a reference consensus direction (e.g., the Yajie 
center or a ground-truth label), define for each expert $E_i$ the 
\textbf{deviation 0-form} $f_i \in \Omega^0(\cM)$:
\begin{equation}
  f_i(x) = \|E_i(x) - \bar{E}(x)\|_{\text{signed}},
  \label{eq:deviation}
\end{equation}
where $\bar{E}(x)$ is the consensus output at $x$ and $\|\cdot\|_{\text{signed}}$ 
is a signed distance (positive for deviation in one direction, negative for 
the opposite). The \textbf{audit 1-form} is then:
\begin{equation}
  A = d_0\left(\frac{1}{M}\sum_{i=1}^{M} f_i\right) \in \Omega^1(\cM),
  \label{eq:audit1form}
\end{equation}
representing the spatial gradient of the average expert deviation. In the 
discrete setting, $A$ is a 1-cochain on the simplicial complex.
\end{definition}

\bigskip
\noindent
\textit{\textbf{define1-} $M$ each $E_i$ define\textbf{0-} $f_i(x) = \|E_i(x) - \bar{E}(x)\|_{\text{signed}}$where $\bar{E}(x)$  $x$ \textbf{1-} $A = d_0(M^{-1}\sum_i f_i)$}

% ---------------------------------------------------------------------------
\subsection{Persistent Homology}
\subsection*{}

\begin{definition}[Persistence Module / ]
\label{def:pers}
Let $\{\cM_{\varepsilon}\}_{\varepsilon \geq 0}$ be the density sublevel 
filtration. For each $k \geq 0$, the $k$-th homology functor $H_k(\cdot; \mathbb{F})$ 
(with coefficients in a field $\mathbb{F}$, typically $\mathbb{Z}_2$ or $\mathbb{Q}$) 
applied to the filtration yields a \textbf{persistence module} 
$\mathbb{V}_k = \{H_k(\cM_{\varepsilon}; \mathbb{F})\}_{\varepsilon \geq 0}$ 
with linear maps $v_{\varepsilon_1, \varepsilon_2}: H_k(\cM_{\varepsilon_1}; \mathbb{F}) 
\to H_k(\cM_{\varepsilon_2}; \mathbb{F})$ induced by inclusions 
$\cM_{\varepsilon_1} \hookrightarrow \cM_{\varepsilon_2}$ for 
$\varepsilon_1 \leq \varepsilon_2$.
\end{definition}

\bigskip
\noindent
\textit{\textbf{define \ref{def:pers}} $\{\cM_{\varepsilon}\}_{\varepsilon \geq 0}$ each $k \geq 0$ $k$  $H_k(\cdot; \mathbb{F})$ $\mathbb{F}$ \textbf{} $\mathbb{V}_k = \{H_k(\cM_{\varepsilon}; \mathbb{F})\}_{\varepsilon \geq 0}$}

By the structure theorem of persistent homology \cite{zomorodian2005, edelsbrunner2008}, 
$\mathbb{V}_k$ decomposes into a direct sum of \textbf{interval modules}:
\begin{equation}
  \mathbb{V}_k \cong \bigoplus_{j \in \mathcal{J}_k} \mathbb{I}_{[b_j, d_j]},
  \label{eq:interval_decomp}
\end{equation}
where each $\mathbb{I}_{[b, d]}$ is an indecomposable module supported on 
the interval $[b, d]$, with \textbf{birth} $b_j$ (the density threshold at 
which the homology class appears) and \textbf{death} $d_j$ (the threshold 
at which it merges with another class or becomes trivial).

\bigskip
\noindent
\textit{theorem$\mathbb{V}_k$ \textbf{}each $\mathbb{I}_{[b, d]}$  $[b, d]$ \textbf{} $b_j$\textbf{} $d_j$}

\begin{definition}[Persistence Diagram / ]
The \textbf{$k$-dimensional persistence diagram} is the multiset:
\begin{equation}
  \cD_k = \{(b_j, d_j) \in \R^2 : b_j < d_j, \; j \in \mathcal{J}_k\},
  \label{eq:pers_diagram}
\end{equation}
together with all points on the diagonal $\{(b, b) : b \geq 0\}$ taken with 
infinite multiplicity (representing classes of zero persistence).

The \textbf{persistence} (or \textbf{lifetime}) of a homology class is 
$\pers([\gamma]) = d_j - b_j$. Classes with large persistence are 
considered \textbf{topologically significant}, while those with small 
persistence are typically attributed to noise \cite{chazal2017}.
\end{definition}

\bigskip
\noindent
\textit{\textbf{define} $k$ \textbf{} $\cD_k = \{(b_j, d_j) : b_j < d_j\}$all\textbf{}\textbf{} $\pers([\gamma]) = d_j - b_j$\textbf{}}

% ===========================================================================
%  SECTION 3: AUDIT INSTANTONS — FORMAL DEFINITION
% ===========================================================================
\section{Audit Instantons: Formal Definition}
\section*{Instantondefine}

\subsection{The Audit Vector Field and Its Cohomology}
\subsection*{}

Consider an audit pool of $M$ experts on the Situs manifold $\cM$. Let 
$\{x_{\alpha}\}_{\alpha=1}^{N} \subset \cM$ be a set of $N$ audit data 
points (which may or may not coincide with training data). At each point 
$x$, expert $E_i$ produces a prediction $E_i(x) \in \R^d$.

\begin{definition}[Expert Deviation Field / ]
\label{def:deviation_field}
The \textbf{signed expert deviation field} is a function 
$\mathbf{f}: \cM \to \R^M$ defined componentwise:
\begin{equation}
  \mathbf{f}(x) = (f_1(x), f_2(x), \ldots, f_M(x)),
  \label{eq:f_vector}
\end{equation}
where each $f_i: \cM \to \R$ is the signed deviation of expert $i$ from 
a reference value $y(x)$ (e.g., ground truth or Yajie consensus):
\begin{equation}
  f_i(x) = \langle E_i(x) - y(x), \mathbf{n}(x) \rangle,
  \label{eq:signed_dev}
\end{equation}
with $\mathbf{n}(x)$ a distinguished normal direction field (for scalar 
outputs, $\mathbf{n}(x) \equiv 1$ and $f_i(x) = E_i(x) - y(x)$).
\end{definition}

\bigskip
\noindent
\textit{\textbf{define \ref{def:deviation_field}}\textbf{} $\mathbf{f}: \cM \to \R^M$whereeach $f_i(x) = \langle E_i(x) - y(x), \mathbf{n}(x) \rangle$  $i$  $y(x)$ }

\begin{definition}[Audit 1-Form / 1-]
\label{def:A}
The \textbf{audit 1-form} $A \in \Omega^1(\cM)$ is the exterior derivative 
of the \emph{mean} signed deviation:
\begin{equation}
  A = d_0 \bar{f}, \quad \text{where} \quad 
  \bar{f}(x) = \frac{1}{M} \sum_{i=1}^{M} f_i(x).
  \label{eq:A_def}
\end{equation}
Equivalently, in local coordinates $(x^1, \ldots, x^n)$:
\begin{equation}
  A = \sum_{\mu=1}^{n} \frac{\partial \bar{f}}{\partial x^{\mu}} \, dx^{\mu}.
  \label{eq:A_local}
\end{equation}
\end{definition}

\bigskip
\noindent
\textit{\textbf{define \ref{def:A}1-}\textbf{1-} $A \in \Omega^1(\cM)$ $A = d_0 \bar{f}$}

\begin{proposition}[Automatic Closedness / ]
\label{prop:closed}
By construction, $d_1 A = d_1 d_0 \bar{f} = 0$. Hence $A$ is always a 
\emph{closed} 1-form: $A \in \ker(d_1) = Z^1(\cM)$.
\end{proposition}

\bigskip
\noindent
\textit{\textbf{Proposition \ref{prop:closed}}$d_1 A = d_1 d_0 \bar{f} = 0$therefore $A$ \textbf{}1-$A \in Z^1(\cM)$}

\begin{proposition}[Exactness Condition / ]
\label{prop:exactness}
The audit 1-form $A$ is \emph{exact} (i.e., $A \in \im(d_0) = B^1(\cM)$) 
if and only if the mean deviation $\bar{f}$ is well-defined as a global 
smooth function. In that case, $A$ lies in the trivial cohomology class: 
$[A] = 0 \in H^1_{\text{dR}}(\cM)$.

If $\bar{f}$ has \emph{monodromy} — i.e., its integral around a closed loop 
$\gamma$ is non-zero:
\begin{equation}
  \oint_{\gamma} A = \oint_{\gamma} d_0 \bar{f} \neq 0,
  \label{eq:monodromy}
\end{equation}
then $\bar{f}$ is not globally single-valued, and $A$ represents a 
\emph{non-trivial} cohomology class $[A] \neq 0 \in H^1_{\text{dR}}(\cM)$.
\end{proposition}

\bigskip
\noindent
\textit{\textbf{Proposition \ref{prop:exactness}}1- $A$ \textbf{}i.e. $A \in B^1(\cM)$if and only if $\bar{f}$ define $A$  $\bar{f}$ \textbf{}—— $\gamma$ —— $\bar{f}$ $A$ \textbf{}}

% ---------------------------------------------------------------------------
\subsection{Definition of Audit Instanton}
\subsection*{Instantondefine}

\begin{definition}[Audit Instanton / Instanton]
\label{def:instanton}
Let $(\cM, \rho, \cE)$ be a Situs manifold and $A \in \Omega^1(\cM)$ the 
audit 1-form. An \textbf{audit instanton} is a piecewise-smooth closed curve 
$\gamma: S^1 \to \cM$ (a 1-cycle) satisfying the following three conditions:

\begin{enumerate}[label=\textbf{(AI\arabic*)}]
  \item \textbf{Low-Density Localization / }: 
        $\gamma$ is contained in a low-density region:
        \begin{equation}
          \gamma \subset \cM_{\rho_{\text{crit}}} = \{x \in \cM : \rho(x) < \rho_{\text{crit}}\}.
          \label{eq:AI1}
        \end{equation}
  
  \item \textbf{Local Consistency / Consistency}: 
        The audit 1-form is locally exact along $\gamma$:
        \begin{equation}
          \forall p \in \gamma, \; \exists \text{ a neighborhood } U_p \subset \cM 
          \text{ and } g_p \in \Omega^0(U_p) \text{ such that } A|_{U_p} = d_0 g_p.
          \label{eq:AI2}
        \end{equation}
        This holds automatically by Proposition \ref{prop:exactness} — the 
        substantive content is that the local primitive $g_p$ must have 
        \emph{bounded variation} along $\gamma$: 
        $\|\nabla g_p\|_{L^{\infty}(\gamma \cap U_p)} \leq C$ for some 
        universal constant $C$.
  
  \item \textbf{Global Non-Triviality / }: 
        The cycle is not a boundary in the low-density region:
        \begin{equation}
          \gamma \notin \im(\partial_2) \text{ in } H_1(\cM_{\rho_{\text{crit}}}; \Z),
          \label{eq:AI3}
        \end{equation}
        where $\partial_2: C_2(\cM_{\rho_{\text{crit}}}) \to C_1(\cM_{\rho_{\text{crit}}})$ 
        is the simplicial boundary operator. Equivalently, $\gamma$ represents 
        a \emph{non-trivial} homology class:
        $[\gamma] \neq 0 \in H_1(\cM_{\rho_{\text{crit}}}; \Z)$.

        Moreover, the \textbf{holonomy} (circulation) of $A$ around $\gamma$ 
        is non-zero:
        \begin{equation}
          \Phi(\gamma) = \oint_{\gamma} A > \eta,
          \label{eq:holonomy}
        \end{equation}
        where $\eta > 0$ is a significance threshold (e.g., the expected 
        noise level of the audit 1-form).
  \end{enumerate}
\end{definition}

\bigskip
\noindent
\textit{\textbf{define \ref{def:instanton}Instanton} $(\cM, \rho, \cE)$ Situs$A$ 1-\textbf{Instanton}satisfying $\gamma: S^1 \to \cM$1-\textbf{(AI1) }$\gamma \subset \cM_{\rho_{\text{crit}}}$\textbf{(AI2) Consistency}$A$  $\gamma$  $\gamma$ \textbf{(AI3) }$\gamma \notin \im(\partial_2)$  $H_1(\cM_{\rho_{\text{crit}}}; \Z)$  $A$  $\gamma$ \textbf{}$\Phi(\gamma) = \oint_{\gamma} A > \eta$}

\begin{remark}[Physical Interpretation / ]
The three conditions capture the essential ``instanton-ness'' of the 
configuration:
\begin{itemize}
  \item (AI1) ensures we are in the ``weak-coupling'' (low-data) regime 
        where non-perturbative effects dominate — analogous to instantons 
        contributing at small $g$ in QFT.
  \item (AI2) ensures that \emph{locally} the deviation field looks like 
        a gradient of a well-behaved function — analogous to the classical 
        instanton solving the equations of motion. Pointwise Yajie audit 
        sees local consistency and concludes ``high consensus.''
  \item (AI3) ensures that \emph{globally} the deviation accumulates 
        non-zero net error — the topological obstruction that pointwise 
        audit cannot detect. This is the ``instantonic charge'' or 
        ``winding number.''
\end{itemize}
\end{remark}

\bigskip
\noindent
\textit{\textbf{remark}``Instanton''(AI1) ``''(AI2) \textbf{}——YajieConsistencyConclusion``''(AI3) \textbf{}——i.e.``Instanton''``''}

% ---------------------------------------------------------------------------
\subsection{Non-Perturbative Character}
\subsection*{}

The defining property of audit instantons — and what makes them 
potentially catastrophic — is their \textbf{invisibility to pointwise 
audit}. We formalize this as follows.

\begin{proposition}[Pointwise Audit Blindness / ]
\label{prop:blindness}
Let $\gamma$ be an audit instanton. The standard pointwise Yajie audit 
score $\Yajie(x)$, defined as a local function of expert outputs at $x$ 
alone, satisfies:
\begin{equation}
  \inf_{x \in \gamma} \Yajie(x) \geq \Yajie_{\text{threshold}},
  \label{eq:blindness}
\end{equation}
while the \emph{ground-truth error} integrated along $\gamma$ exceeds the 
audit tolerance:
\begin{equation}
  \frac{1}{|\gamma|} \oint_{\gamma} |\bar{f}(x)| \, dx \gg \varepsilon_{\text{tol}}.
  \label{eq:error}
\end{equation}
\end{proposition}

\bigskip
\noindent
\textit{\textbf{Proposition \ref{prop:blindness}} $\gamma$ InstantonYajieScore $\Yajie(x)$ satisfying $\inf_{x \in \gamma} \Yajie(x) \geq \Yajie_{\text{threshold}}$ $\gamma$ \textbf{}}

\begin{proof}[Proof Sketch / Proof Sketch]
Condition (AI2) guarantees that at each point $x \in \gamma$, the local 
deviation has bounded variation — all experts deviate ``smoothly'' in the 
same direction. Pointwise Yajie, which only compares expert outputs at $x$, 
sees this as consensus (everyone agrees) and assigns a high score. However, 
condition (AI3) ensures that the accumulated deviation around the full 
cycle $\gamma$ is non-zero — the consensus is \emph{spurious} and the 
total error is systematic. This is a direct analog of the perturbative 
expansion's inability to detect $\sim \exp(-1/g^2)$ terms.
\end{proof}

\bigskip
\noindent
\textit{\textbf{Proof Sketch}(AI2) $\gamma$  $x$——all``''Yajie $x$ allhowever(AI3) $\gamma$ ——\textbf{} $\sim \exp(-1/g^2)$ }

% ===========================================================================
%  SECTION 4: DETECTION THEORY
% ===========================================================================
\section{Detection Theory via Persistent Homology}
\section*{}

\subsection{Persistent Homology of the Density Filtration}
\subsection*{}

The density sublevel filtration $\{\cM_{\varepsilon}\}_{\varepsilon \geq 0}$ 
captures the topology of data-sparse regions at multiple scales. As $\varepsilon$ 
increases, new topological features (connected components, cycles, voids) are 
born, persist for some $\varepsilon$-range, and eventually die.

\begin{definition}[Density Filtration Barcode / ]
\label{def:barcode}
The \textbf{1-dimensional barcode} of the density filtration is the set of 
intervals:
\begin{equation}
  \cB_1 = \{(b_j, d_j)\}_{j \in \mathcal{J}_1},
  \label{eq:barcode}
\end{equation}
where each interval corresponds to a 1-cycle (``hole'' or ``tunnel'') that:
\begin{itemize}
  \item Is \textbf{born} at density $b_j$: the 1-cycle first appears as a 
        hole in $\cM_{b_j}$;
  \item \textbf{Dies} at density $d_j$: the hole is filled (becomes a 
        boundary) in $\cM_{d_j}$.
\end{itemize}
Equivalently, this is the persistence diagram $\cD_1$ of the superlevel 
set filtration of $\rho$, viewed in the reverse parametrization.
\end{definition}

\bigskip
\noindent
\textit{\textbf{define \ref{def:barcode}}\textbf{1} $\cB_1 = \{(b_j, d_j)\}_{j \in \mathcal{J}_1}$each1-``''``'' $b_j$ \textbf{}$\cM_{b_j}$  $d_j$ \textbf{}$\cM_{d_j}$ }

The key intuition is that \textbf{long-lived 1-cycles} (those with large 
$\pers = d_j - b_j$) correspond to \emph{robust} low-density 
regions — tunnels or loops in the data manifold that remain empty of 
data over a wide range of density thresholds. These are precisely the 
regions where audit instantons can form.

\bigskip
\noindent
\textit{\textbf{1-} $\pers = d_j - b_j$ \textbf{}——Instanton}

% ---------------------------------------------------------------------------
\subsection{Correspondence: 1-Cycles and Audit Instanton Candidates}
\subsection*{1-Instanton}

We now establish the connection between persistent 1-cycles and audit instantons.

\begin{definition}[Audit Instanton Candidate / Instanton]
\label{def:candidate}
A \textbf{candidate audit instanton} is a triple $(\gamma, b, d)$, where:
\begin{enumerate}[label=(\roman*)]
  \item $\gamma \subset \cM$ is a closed 1-cycle representing a class 
        $[\gamma] \in H_1(\cM_{\varepsilon}; \Z)$;
  \item $(b, d)$ is the persistence interval associated to $[\gamma]$, 
        with birth $b$ and death $d$;
  \item $d > \rho_{\text{crit}}$ (the cycle persists beyond the critical 
        density threshold), indicating the cycle is \textbf{not a noise artifact}.
\end{enumerate}
The candidate is \textbf{charged} if additionally $\oint_{\gamma} A > \eta$, 
i.e., the audit 1-form has non-trivial holonomy along $\gamma$.
\end{definition}

\bigskip
\noindent
\textit{\textbf{define \ref{def:candidate}Instanton}\textbf{Instanton} $(\gamma, b, d)$where(i) $\gamma \subset \cM$ 1- $H_1(\cM_{\varepsilon}; \Z)$ (ii) $(b, d)$  $[\gamma]$ (iii) $d > \rho_{\text{crit}}$ $\oint_{\gamma} A > \eta$\textbf{}}

\begin{proposition}[Filtration-Persistence Correspondence / -]
\label{prop:filtration}
If $\gamma$ is an audit instanton in the sense of Definition \ref{def:instanton}, 
then $\gamma$ corresponds to a 1-cycle in the density filtration with:
\begin{equation}
  \Birth(\gamma) \leq \rho_{\text{crit}} \quad \text{and} \quad 
  \Death(\gamma) \geq \max_{x \in \text{interior}(\gamma)} \rho(x),
  \label{eq:birthdeath}
\end{equation}
where $\text{interior}(\gamma)$ denotes the minimal spanning surface (2-chain) 
bounded by $\gamma$ in $\cM$.

In particular, the persistence satisfies 
$\pers(\gamma) \geq \max_{x \in \text{interior}(\gamma)} \rho(x) - \rho_{\text{crit}}$.
\end{proposition}

\bigskip
\noindent
\textit{\textbf{Proposition \ref{prop:filtration}-} $\gamma$ define \ref{def:instanton} Instanton $\gamma$ 1- $\leq \rho_{\text{crit}}$ $\geq \max_{x \in \text{interior}(\gamma)} \rho(x)$in particularsatisfying $\pers(\gamma) \geq \max_{x \in \text{interior}(\gamma)} \rho(x) - \rho_{\text{crit}}$}

% ---------------------------------------------------------------------------
\subsection{The Main Structural Theorem}
\subsection*{theorem}

We now state and prove the central result of this paper.

\begin{theorem}[Existence of Audit Instanton Candidates / Instantonexistence]
\label{thm:main}
Let $(\cM, \rho, \cE)$ be a Situs manifold with the following properties:
\begin{enumerate}[label=(H\arabic*)]
  \item \textbf{Smoothness / }: $\rho \in C^2(\cM, \R_{\geq 0})$ and 
        $\cM$ has non-empty 1-homology: $\rank H_1(\cM; \Z) \geq 1$.
  \item \textbf{Bounded Expert Deviation / }: There exists 
        a constant $K > 0$ such that $\|\nabla \bar{f}\|_{L^{\infty}(\cM)} < K$, 
        where $\bar{f}$ is the mean signed deviation.
  \item \textbf{Non-Degeneracy / }: The critical density satisfies 
        $\rho_{\text{crit}} > 0$ and the sublevel set 
        $\cM_{\rho_{\text{crit}}}$ is non-empty and has non-trivial 
        1-homology (at least one ``hole'' survives).
  \item \textbf{Sufficient Expert Pool / }: $M \geq 2$, and 
        expert deviations are not perfectly anti-correlated (to exclude 
        the trivial case where all deviations cancel exactly).
\end{enumerate}

Then every non-trivial 1-homology class 
$0 \neq [\alpha] \in H_1(\cM_{\rho_{\text{crit}}}; \Z)$ contains at least 
one \textbf{candidate audit instanton}: there exists a 1-cycle 
$\gamma$ representing $[\alpha]$ such that $(\gamma, b, d)$ is a candidate 
in the sense of Definition \ref{def:candidate}.

Moreover, if the expert deviation field satisfies the \emph{positivity 
condition}:
\begin{equation}
  \forall i \in \{1,\ldots,M\}, \; \forall x \in \cM_{\rho_{\text{crit}}}:
  \quad f_i(x) \cdot \bar{f}(x) \geq 0,
  \label{eq:positivity}
\end{equation}
(i.e., all experts have the same sign of deviation on low-density regions), 
then the candidate is \textbf{charged}: $\oint_{\gamma} A > \eta$ for some 
$\eta > 0$ depending only on $\rho_{\text{crit}}$ and $K$.
\end{theorem}

\bigskip
\noindent
\textit{\textbf{theorem \ref{thm:main}Instantonexistence} $(\cM, \rho, \cE)$ satisfyingSitus(H1) $\rho \in C^2$  $\rank H_1(\cM; \Z) \geq 1$(H2) $\|\nabla \bar{f}\|_{\infty} < K$(H3) $\rho_{\text{crit}} > 0$  $\cM_{\rho_{\text{crit}}}$ 1-(H4) $M \geq 2$  $\cM_{\rho_{\text{crit}}}$ each1-\textbf{Instanton}furthermoresatisfyingall\textbf{}}

\begin{proof}[Proof / Proof]
We proceed in three steps.

\medskip
\noindent
\textit{}
\textbf{Step 1. Representation of the homology class.}
Let $0 \neq [\alpha] \in H_1(\cM_{\rho_{\text{crit}}}; \Z)$. By standard 
results in differential topology \cite{bott_tu}, there exists a smooth 
embedded closed curve $\gamma_0: S^1 \to \cM_{\rho_{\text{crit}}}$ 
representing $[\alpha]$. Since $\cM_{\rho_{\text{crit}}}$ is an open 
submanifold of $\cM$ (as $\rho$ is continuous), we may perturb $\gamma_0$ 
slightly to ensure transversality with respect to the foliation defined 
by the level sets of $\rho$.

\bigskip
\noindent
\textit{}
\textbf{Step 2. Persistence in the density filtration.}
Consider the image of $[\alpha]$ under the inclusion 
$\iota_{\varepsilon}: \cM_{\rho_{\text{crit}}} \hookrightarrow \cM_{\varepsilon}$ 
for $\varepsilon \geq \rho_{\text{crit}}$. The class $[\alpha]$ is born at 
some density $b \leq \rho_{\text{crit}}$ (by definition, it exists in 
$\cM_{\rho_{\text{crit}}}$) and dies at some $d \geq \rho_{\text{crit}}$ 
when it becomes a boundary. Since $\rank H_1(\cM_{\varepsilon}; \Z)$ is 
non-decreasing in $\varepsilon$ (by functoriality of homology) and 
$[\alpha]$ is non-trivial in $\cM_{\rho_{\text{crit}}}$, its death must 
satisfy $d > \rho_{\text{crit}}$. Thus $\gamma_0$ corresponds to a 
1-cycle in the persistence barcode with $\pers(\gamma_0) = d - b > 0$, 
satisfying condition (iii) of Definition \ref{def:candidate}.

\bigskip
\noindent
\textit{}
\textbf{Step 3. Charging condition.}
We now verify the holonomy condition. Under hypothesis (H2), the audit 
1-form $A = d_0 \bar{f}$ satisfies $\|A\|_{L^{\infty}(\cM)} < K$. For any 
1-cycle $\gamma$,
\begin{equation}
  \left|\oint_{\gamma} A\right| \leq K \cdot \length(\gamma).
  \label{eq:hol_bound}
\end{equation}

Under the positivity condition (\ref{eq:positivity}), for any $x$ in 
$\cM_{\rho_{\text{crit}}}$, all $f_i(x)$ have the same sign as $\bar{f}(x)$. 
This implies that on $\cM_{\rho_{\text{crit}}}$, the absolute mean deviation 
$|\bar{f}|$ is bounded below by a positive constant 
$\delta = \min_{x \in \cM_{\rho_{\text{crit}}}} |\bar{f}(x)| > 0$ (using 
compactness of $\cM$ and continuity of $\bar{f}$). 

Now, consider a refined representative $\gamma$ of $[\alpha]$ that 
minimizes the length-to-holonomy ratio among all representatives. By 
the isoperimetric inequality on $\cM$, there exists a constant 
$C_{\cM} > 0$ such that for any non-trivial 1-cycle $\gamma$ representing 
$[\alpha]$:
\begin{equation}
  \left|\oint_{\gamma} A\right| \geq \frac{\delta \cdot \inf_{\Sigma: \partial\Sigma = \gamma} \vol(\Sigma)}{C_{\cM} \cdot \length(\gamma)},
  \label{eq:isoperimetric}
\end{equation}
where $\Sigma$ is any 2-chain bounding $\gamma$ in $\cM$.

Since $[\alpha] \neq 0$ in $H_1(\cM_{\rho_{\text{crit}}}; \Z)$, any surface 
$\Sigma$ bounding $\gamma$ must intersect the complement 
$\cM \setminus \cM_{\rho_{\text{crit}}}$ (i.e., high-density regions), and 
therefore has volume at least $\vol(\cM_{\rho_{\text{crit}}}) > 0$. Hence:
\begin{equation}
  \left|\oint_{\gamma} A\right| \geq \eta := 
    \frac{\delta \cdot \vol(\cM_{\rho_{\text{crit}}})}{C_{\cM} \cdot \diam(\cM)} > 0.
  \label{eq:eta_bound}
\end{equation}
This establishes condition (AI3) of Definition \ref{def:instanton} with 
the explicit threshold $\eta$ given above.

Finally, condition (AI2) (local exactness) follows from Proposition 
\ref{prop:exactness} and the smoothness of $\bar{f}$ on the compact set 
$\gamma$, which guarantees bounded variation of the local primitive.

\medskip
\noindent
\textit{$\gamma$ define \ref{def:candidate} satisfyingtheoremall}
We have thus constructed, for every non-trivial $[\alpha] \in H_1(\cM_{\rho_{\text{crit}}}; \Z)$, 
a candidate audit instanton $\gamma$ representing $[\alpha]$, completing the proof.
\end{proof}

\begin{corollary}[Minimal Number of Audit Instantons / Instanton]
\label{cor:min}
Under hypotheses (H1)--(H4), the minimum number of \emph{distinct} audit 
instanton candidates in $(\cM, \rho, \cE)$ is at least 
$\rank H_1(\cM_{\rho_{\text{crit}}}; \Z)$ — the first Betti number of the 
low-density region.
\end{corollary}

\bigskip
\noindent
\textit{\textbf{corollary \ref{cor:min}Instanton}assumption (H1)--(H4) $(\cM, \rho, \cE)$ \textbf{}Instanton $\rank H_1(\cM_{\rho_{\text{crit}}}; \Z)$——Betti}

% ===========================================================================
%  SECTION 5: RELATION TO YAJIE CONSENSUS
% ===========================================================================
\section{Invisibility to Yajie Consensus Scoring}
\section*{Yajie}

The Yajie consensus score $\Yajie(x)$ at a point $x \in \cM$ is a local 
functional of the expert outputs $\{E_i(x)\}_{i=1}^{M}$ that measures the 
degree of agreement among experts. It is typically defined as (the negative 
of) the variance, or via the discrete Hodge norm of the deviation field 
\cite{scx_framework}:
\begin{equation}
  \Yajie(x) = 1 - \frac{\operatorname{Var}_i(f_i(x))}{\operatorname{Var}_{\max}},
  \label{eq:yajie_def}
\end{equation}
where $\operatorname{Var}_{\max}$ is a normalization constant.

\bigskip
\noindent
\textit{YajieScore $\Yajie(x)$ Consistencydefine}

\begin{theorem}[Yajie Blindness to Audit Instantons / YajieInstanton]
\label{thm:yajie_blindness}
Let $\gamma$ be a charged audit instanton satisfying Definition 
\ref{def:instanton} with positivity condition (\ref{eq:positivity}). 
Then:
\begin{enumerate}[label=(\roman*)]
  \item \textbf{Pointwise consensus is high / }: 
        There exists $\theta > 0$ such that for all $x \in \gamma$:
        \begin{equation}
          \Yajie(x) \geq 1 - \frac{K^2 \cdot (\diam(U_x))^2}{M \cdot \operatorname{Var}_{\max}} > 1 - \theta,
          \label{eq:yajie_high}
        \end{equation}
        where $U_x$ is the local neighborhood of $x$ on which $A$ is exact.
  
  \item \textbf{Cycle-integrated error is high / }: 
        The integrated deviation around $\gamma$ is bounded below:
        \begin{equation}
          \frac{1}{|\gamma|} \oint_{\gamma} |\bar{f}(x)| \, dx \geq 
          \frac{|\Phi(\gamma)|}{|\gamma|} > \frac{\eta}{|\gamma|},
          \label{eq:cycle_error}
        \end{equation}
        where $\Phi(\gamma) = \oint_{\gamma} A$ is the holonomy.
  
  \item \textbf{Ratio diverges with persistence / }: 
        As $\pers(\gamma) \to \infty$, the ratio
        \begin{equation}
          \frac{\text{Pointwise Yajie estimate of audit quality}}
               {\text{Cycle-integrated ground-truth audit quality}} \to \infty.
          \label{eq:divergence}
        \end{equation}
\end{enumerate}
\end{theorem}

\bigskip
\noindent
\textit{\textbf{theorem \ref{thm:yajie_blindness}YajieInstanton} $\gamma$ satisfyingdefine \ref{def:instanton} Instanton(i) Yajie $\gamma$ Score(ii)  $\gamma$ lower bound(iii) Growth}

\begin{proof}
For (i): By (AI2), on each local neighborhood $U_x$, $A = d_0 g_x$ with 
$\|\nabla g_x\|_{L^{\infty}} \leq C$. Hence the variation of $\bar{f}$
across $U_x$ is bounded by $C \cdot \diam(U_x)$. The variance of 
$\{f_i\}$ on $U_x$ is then bounded by $O(\diam(U_x)^2)$, since all experts 
deviate similarly on the low-density region (positivity condition). 
The Yajie score approaches 1 as $\diam(U_x) \to 0$ or as $M \to \infty$.

For (ii): By the holonomy condition (AI3), 
$\oint_{\gamma} A = \oint_{\gamma} d_0 \bar{f}$ is non-zero. Since 
$\bar{f}$ is continuous on the compact curve $\gamma$, there exists 
a lower bound on $|\bar{f}|$ proportional to $|\Phi(\gamma)|/|\gamma|$.

For (iii): As $\pers(\gamma) \to \infty$, the cycle $\gamma$ wraps around 
a region where $\rho$ remains arbitrarily small over a large range. In 
this regime, the local Yajie score is maximized (all experts agree due to 
shared extrapolation bias), while the integrated error is minimized only 
down to the topological lower bound — achieving a finite positive floor. 
The ratio thus diverges.
\end{proof}

\bigskip
\noindent
\textit{\textbf{Proof}(i) (AI2)each $U_x$  $A = d_0 g_x$$\|\nabla g_x\|_{\infty} \leq C$ $\bar{f}$ allYajieScore1(ii) (AI3)$\oint_{\gamma} A \neq 0$ $\gamma$  $|\bar{f}|$ lower bound(iii)  $\pers(\gamma) \to \infty$$\gamma$  $\rho$ anyYajielower bound}

\begin{remark}[The Non-Perturbative Analogy, Revisited / ]
Theorem \ref{thm:yajie_blindness} is the precise analog of the statement 
that instanton contributions $\sim \exp(-1/g^2)$ are invisible to 
perturbation theory: the pointwise Yajie audit (the ``perturbative expansion'') 
assigns high scores because each local patch looks consistent, yet the 
global topological charge $\Phi(\gamma)$ (the ``instanton number'') 
reveals systematic error.

The key difference from a generic correlated failure is the \emph{topological 
robustness}: an audit instanton cannot be ``audited away'' by adding more 
data points in the vicinity — it requires filling the topological hole, 
i.e., adding data points \emph{inside} the cycle, which may be 
geometrically or physically impossible in many SCX applications.
\end{remark}

\bigskip
\noindent
\textit{\textbf{remark}theorem \ref{thm:yajie_blindness} Instanton $\sim \exp(-1/g^2)$ \textbf{robustness}Instanton``''——i.e.\textbf{}SCX}

% ===========================================================================
%  SECTION 6: DETECTION PROTOCOL
% ===========================================================================
\section{Detection Protocol}
\section*{}

We now provide a concrete, algorithmically implementable protocol for 
detecting audit instantons in real SCX deployments. The protocol assumes 
a discretized Situs manifold represented as a simplicial complex $K$ 
with $N$ vertices (data points).

\bigskip
\noindent
\textit{SCXInstantonassumptionSitus $N$  $K$ }

% ---------------------------------------------------------------------------
\begin{protocol}[Audit Instanton Detection (AID) / Instanton]
\label{prot:aid}
\textbf{Input / }:
\begin{itemize}
  \item Simplicial complex $K$ approximating $\cM$, with vertex set 
        $V = \{v_1, \ldots, v_N\}$;
  \item Estimated data density $\hat{\rho}: V \to \R_{\geq 0}$ (e.g., 
        kernel density estimate);
  \item Expert outputs $\{E_i(v)\}_{i=1}^{M}$ for all $v \in V$;
  \item Critical density threshold $\rho_{\text{crit}}$;
  \item Persistence significance threshold $\tau > 0$;
  \item Holonomy (circulation) threshold $\eta > 0$.
\end{itemize}

\bigskip
\noindent
\textit{\textbf{ \ref{prot:aid}Instanton}}
\textit{\textbf{} $K$  $\cM$ $\hat{\rho}$ $\{E_i\}$ $\rho_{\text{crit}}$ $\tau$ $\eta$}

\medskip
\noindent
\textbf{Output / }: 
A set $\cI$ of detected audit instanton candidates, each specified as a 
1-cycle with associated metadata (persistence, holonomy, consensus score, 
cycle geometry).

\bigskip
\noindent
\textit{\textbf{}Instanton $\cI$each1-Score}

\medskip
\noindent
\textbf{Step 1: Density Estimation and Filtration / 1}
\begin{enumerate}[label=\arabic*.]
  \item Estimate $\hat{\rho}: V \to \R_{\geq 0}$ using kernel density 
        estimation (KDE) with bandwidth $h$ chosen by cross-validation:
        \begin{equation}
          \hat{\rho}(v) = \frac{1}{N h^n} \sum_{j=1}^{N} 
            K\left(\frac{\|v - v_j\|}{h}\right).
          \label{eq:kde}
        \end{equation}
  \item Extend $\hat{\rho}$ to $K$ via linear interpolation on simplices.
  \item Construct the \textbf{density sublevel filtration}: 
        for each vertex $v$, assign filtration value $\hat{\rho}(v)$; 
        for each edge/simplex, assign the maximum density of its vertices 
        (lower-star filtration).

        \textit{each $v$  $\hat{\rho}(v)$/}
\end{enumerate}

\textbf{Step 2: Persistent Homology Computation / 2}
\begin{enumerate}[label=\arabic*.]
  \setcounter{enumi}{2}
  \item Compute the 1-dimensional persistent homology 
        $\PH_1(K, \hat{\rho})$ using the standard matrix reduction 
        algorithm \cite{edelsbrunner2008} or Ripser \cite{bauer2021}.
  
        \textit{Ripser1 $\PH_1(K, \hat{\rho})$}

  \item Extract the persistence diagram $\cD_1 = \{(b_j, d_j)\}_{j}$ and 
        identify \textbf{significant 1-cycles}: those with persistence 
        $\pers_j = d_j - b_j > \tau$.
        
        \textit{ $\cD_1$ \textbf{1-} $\pers_j > \tau$ }

  \item For each significant 1-cycle, compute a \textbf{geometric 
        representative} $\gamma_j$ — a closed edge-path in $K$ that 
        represents the homology class. This can be done via the 
        ``representative cycles'' functionality of standard PH libraries.
        
        \textit{each1-\textbf{} $\gamma_j$—— $K$ }
\end{enumerate}

\textbf{Step 3: Audit 1-Form and Holonomy / 31-}
\begin{enumerate}[label=\arabic*.]
  \setcounter{enumi}{5}
  \item Compute the mean signed deviation $\bar{f}: V \to \R$:
        \begin{equation}
          \bar{f}(v) = \frac{1}{M} \sum_{i=1}^{M} f_i(v),
          \quad f_i(v) = \langle E_i(v) - y(v), \mathbf{n}(v) \rangle.
          \label{eq:f_bar_disc}
        \end{equation}
        
        \textit{ $\bar{f}(v) = M^{-1} \sum_i f_i(v)$}
  
  \item Define the \textbf{discrete audit 1-form} $A$ on each oriented 
        edge $e = [u, v]$:
        \begin{equation}
          A(e) = \bar{f}(v) - \bar{f}(u).
          \label{eq:A_disc}
        \end{equation}
        This is precisely $(d_0 \bar{f})(e)$ in discrete Hodge theory.
        
        \textit{define\textbf{1-} $A(e) = \bar{f}(v) - \bar{f}(u)$Hodge $(d_0 \bar{f})(e)$}

  \item For each candidate cycle $\gamma_j = (v_0, v_1, \ldots, v_k = v_0)$, 
        compute the \textbf{holonomy} (discrete circulation):
        \begin{equation}
          \Phi(\gamma_j) = \sum_{\ell=0}^{k-1} A([v_{\ell}, v_{\ell+1}]).
          \label{eq:holonomy_disc}
        \end{equation}
        
        \textit{each $\gamma_j$\textbf{}$\Phi(\gamma_j) = \sum_{\ell} A([v_{\ell}, v_{\ell+1}])$}
\end{enumerate}

\textbf{Step 4: Consensus Score Verification / 4Score}
\begin{enumerate}[label=\arabic*.]
  \setcounter{enumi}{8}
  \item For each cycle $\gamma_j$ with $|\Phi(\gamma_j)| > \eta$, compute 
        the mean and minimum pointwise Yajie score along $\gamma_j$:
        \begin{equation}
          \bar{\Yajie}(\gamma_j) = \frac{1}{|\gamma_j|} \sum_{v \in \gamma_j} 
            \Yajie(v), \quad 
          \Yajie_{\min}(\gamma_j) = \min_{v \in \gamma_j} \Yajie(v).
          \label{eq:yajie_cycle}
        \end{equation}
        
        \textit{satisfying $|\Phi(\gamma_j)| > \eta$ each $\gamma_j$ $\gamma_j$ YajieScore}
  
  \item Classify cycle $\gamma_j$ as a \textbf{charged audit instanton} if:
        \begin{equation}
          \Yajie_{\min}(\gamma_j) > \Yajie_{\text{threshold}} 
          \quad \text{and} \quad 
          |\Phi(\gamma_j)| > \eta.
          \label{eq:classification}
        \end{equation}
        This is the signature: high local consensus PLUS non-trivial 
        global holonomy.
        
        \textit{ $\Yajie_{\min}(\gamma_j) > \Yajie_{\text{threshold}}$  $|\Phi(\gamma_j)| > \eta$ $\gamma_j$ \textbf{Instanton}}
\end{enumerate}

\textbf{Step 5: Independent Verification / 5}
\begin{enumerate}[label=\arabic*.]
  \setcounter{enumi}{10}
  \item For each charged audit instanton $\gamma_j$, flag the region 
        ``interior'' of $\gamma_j$ (the minimal spanning surface 
        $\Sigma_j$ bounded by $\gamma_j$) for \textbf{independent 
        verification}: additional ground-truth labeling, third-party 
        expert review, or higher-fidelity computational methods.
        
        \textit{eachInstanton $\gamma_j$ $\gamma_j$ ``'' $\gamma_j$  $\Sigma_j$\textbf{}}

  \item Report all detected audit instantons with metadata:
        persistence $(b, d)$, holonomy $\Phi$, Yajie scores, cycle 
        geometry, and recommended verification priority (ordered by 
        $|\Phi| \cdot \pers$).
        
        \textit{allInstanton $(b, d)$ $\Phi$YajieScore $|\Phi| \cdot \pers$ }
\end{enumerate}
\end{protocol}

% ---------------------------------------------------------------------------
\subsection{Computational Complexity}
\subsection*{}

The dominant cost is persistent homology computation, which runs in 
$O(|K|^{\omega})$ time for a simplicial complex $K$, where 
$\omega \approx 2.376$ is the matrix multiplication exponent. In practice, 
for complexes arising from $N \sim 10^3$–$10^4$ data points, computation 
takes seconds to minutes on standard hardware using optimized libraries 
(GUDHI \cite{gudhi}, Ripser \cite{bauer2021}). The remaining steps 
(1-form computation, holonomy, Yajie scores) are $O(N + |E|)$ where $|E|$ 
is the number of edges.

\bigskip
\noindent
\textit{ $O(|K|^{\omega})$for $N \sim 10^3$–$10^4$ GUDHIRipser1-YajieScore $O(N + |E|)$}

% ===========================================================================
%  SECTION 7: EXAMPLES IN SCX
% ===========================================================================
\section{Examples in SCX Auditing Scenarios}
\section*{SCX}

\subsection{Example 1: MLIP for AlN/GaN Materials}
\subsection*{1AlN/GaNMLIP}

Consider the SCX audit of a machine-learned interatomic potential (MLIP) 
for AlN, as described in \cite{scx_framework}. The Situs manifold 
$\cM \subset \R^3$ is parametrized by (bond angle, bond length, coordination 
number). Three experts are audited:

\begin{itemize}
  \item Expert A: trained only on AlN-NaCl phase space;
  \item Expert G: trained only on GaN-wurtzite phase space;
  \item Expert E: trained on full AlGaN phase space.
\end{itemize}

In the region $\Sigma = \{(\theta, r, \kappa) : \theta \in [115^{\circ}, 125^{\circ}], 
r \in [1.75, 1.85] \text{ \AA}\}$ (corresponding to hexagonal graphitic 
coordination), \textbf{no expert has training data}. All three experts 
extrapolate and, due to shared model-class inductive biases, arrive at 
similar (wrong) predictions.

\bigskip
\noindent
\textit{AlNMLIPSCXSitus $\cM \subset \R^3$ AAlN-NaClPhase SpaceGGaN-Phase SpaceEAlGaNPhase Space $\Sigma = \{(\theta, r, \kappa) : \theta \in [115^{\circ}, 125^{\circ}], r \in [1.75, 1.85] \text{ \AA}\}$\textbf{}all}

\begin{example}[The Hexagonal Coordination Instanton / Instanton]
\label{ex:hexagonal}
The data density $\rho$ on $\Sigma$ satisfies $\rho|_{\Sigma} < 10^{-3}$ 
(compared to $\sim 10^{-1}$ in well-sampled regions). The sublevel set 
$\cM_{10^{-3}}$ contains $\Sigma$ and forms a topological annulus 
(1-hole) in the (angle, length) projection — a non-contractible 1-cycle 
$\gamma$ encircling the hexagonal region.

Persistence homology reveals a 1-cycle with $\pers = 0.05$ (birth at 
$\varepsilon = 10^{-3}$, death at $\varepsilon = 0.05$), significantly 
above the noise floor. The holonomy $\Phi(\gamma) = 0.47$ indicates 
systematic deviation around the cycle.

The pointwise Yajie score on $\gamma$ averages 0.89 (high consensus!), 
yet PBE-DFT ground truth shows mean absolute error of 0.31 eV/atom — a 
clear audit instanton.
\end{example}

\bigskip
\noindent
\textit{\textbf{ \ref{ex:hexagonal}Instanton}$\Sigma$  $\rho < 10^{-3}$ $\cM_{10^{-3}}$  $\Sigma$ ——1- $\gamma$  $\pers = 0.05$ $10^{-3}$ $0.05$1- $\Phi(\gamma) = 0.47$ $\gamma$ YajieScore0.89PBE-DFT0.31 eV/atom——Instanton}

% ---------------------------------------------------------------------------
\subsection{Example 2: Adversarial Correlated Failure}
\subsection*{2}

Consider an SCX audit where two experts ($E_1, E_2$) have been adversarially 
manipulated to agree on a specific data submanifold $\Sigma$ known to be 
underrepresented in training. A third honest expert $E_3$ disagrees but 
is outvoted by the Yajie majority rule.

The 1-cycle $\gamma$ bounding $\Sigma$ is an audit instanton: $E_1$ and 
$E_2$ agree (local consistency, AI2), their combined error circulates 
non-trivially around $\gamma$ (global non-triviality, AI3), and the 
region is data-sparse (AI1). Persistent homology detects the hole, 
and the holonomy computation reveals the anomaly even though Yajie 
assigns high consensus based on the 2-to-1 majority.

This illustrates the \textbf{robustness} of audit instanton detection: 
it works even when the majority of experts are systematically wrong 
— precisely the scenario where standard consensus-based audit fails 
most catastrophically.

\bigskip
\noindent
\textit{$E_1, E_2$ $\Sigma$ SCX $E_3$ Yajie $\Sigma$ 1- $\gamma$ Instanton$E_1$  $E_2$ Consistency $\gamma$ ——Yajie21}

% ---------------------------------------------------------------------------
\subsection{Example 3: The SCX ``Information Trap''}
\subsection*{3SCX``''}

A recurring phenomenon in SCX auditing is the \textbf{information trap} 
\cite{scx_framework}: a region of the Situs manifold where all experts 
agree (high Yajie score), no ground truth is available, and yet the 
predictions are systematically wrong. Audit instantons provide a 
topological characterization of information traps: an information trap 
is precisely the interior of a charged audit instanton cycle.

\begin{proposition}[Information Trap = Audit Instanton Interior /  = Instanton]
\label{prop:info_trap}
A connected component $\Omega \subset \cM$ is an \textbf{information trap} 
in the sense of SCX if and only if its boundary $\partial \Omega$ contains 
a charged audit instanton $\gamma$:
\begin{equation}
  \oint_{\gamma} A > \eta, \quad \inf_{x \in \gamma} \Yajie(x) > \Yajie_{\text{threshold}}.
  \label{eq:info_trap_cond}
\end{equation}
\end{proposition}

\bigskip
\noindent
\textit{\textbf{Proposition \ref{prop:info_trap} = Instanton}SCX $\Omega \subset \cM$ \textbf{}if and only if $\partial \Omega$ Instanton $\gamma$}

% ===========================================================================
%  SECTION 8: RELATION TO DISCRETE HODGE THEORY
% ===========================================================================
\section{Relation to Discrete Hodge Theory}
\section*{Hodge}

Audit instantons are naturally situated within the discrete Hodge 
framework of SCX. We establish the precise relationship here.

\subsection{The Hodge Decomposition of the Audit 1-Form}
\subsection*{1-Hodge}

Recall the Hodge decomposition on a compact Riemannian manifold:
\begin{equation}
  \Omega^1(\cM) = \im(d_0) \oplus \mathcal{H}^1(\cM) \oplus \im(\delta_2),
  \label{eq:hodge_decomp}
\end{equation}
where $\mathcal{H}^1(\cM) = \ker(\Delta_1) \cong H^1_{\text{dR}}(\cM)$ is the 
space of harmonic 1-forms. For the audit 1-form $A = d_0 \bar{f}$, we 
clearly have $A \in \im(d_0)$, so the harmonic component vanishes:
\begin{equation}
  A = A_{\text{exact}} + A_{\text{harmonic}} + A_{\text{coexact}} = d_0 \bar{f} + 0 + 0.
  \label{eq:A_decomp}
\end{equation}

\bigskip
\noindent
\textit{RiemannHodge$\Omega^1(\cM) = \im(d_0) \oplus \mathcal{H}^1 \oplus \im(\delta_2)$where $\mathcal{H}^1 \cong H^1_{\text{dR}}$ 1-1- $A = d_0 \bar{f}$ $A \in \im(d_0)$}

\begin{proposition}[Cohomological Signature of Audit Instantons / Instanton]
\label{prop:hodge}
An audit instanon $\gamma$ manifests in discrete Hodge theory as follows:
the cycle $\gamma$ is a \textbf{1-cycle with non-zero evaluation} on the 
harmonic representative of the cohomology class dual to $[\gamma]$. More 
precisely, let $\alpha \in \mathcal{H}^1(\cM)$ be the harmonic 1-form 
Poincar\'{e}-dual to $[\gamma]$. Then:
\begin{equation}
  \oint_{\gamma} \alpha = \int_{\cM} \alpha \wedge \star \eta_{\gamma} \neq 0,
  \label{eq:poincare}
\end{equation}
where $\eta_{\gamma}$ is the Thom form of the normal bundle of $\gamma$.

However, since $A = d_0 \bar{f}$ is exact, its pairing with $\alpha$ is zero:
$\oint_{\gamma} A = 0$ would hold if $\bar{f}$ were globally defined. The 
non-zero holonomy signals that $\bar{f}$ is not globally single-valued — it 
has monodromy around the 1-cycle $\gamma$.
\end{proposition}

\bigskip
\noindent
\textit{\textbf{Proposition \ref{prop:hodge}Instanton}HodgeInstanton $\gamma$  $[\gamma]$ \textbf{1-}howeverbecause $A = d_0 \bar{f}$ $\bar{f}$ ——1- $\gamma$ }

% ---------------------------------------------------------------------------
\subsection{Discrete Laplacian and Audit Instanton Spectrum}
\subsection*{LaplacianInstanton}

\begin{proposition}[Spectral Detection / ]
\label{prop:spectral}
Let $L_1$ be the discrete 1-Laplacian on the simplicial complex $K$ 
approximating $\cM$. The non-trivial 1-cycles corresponding to audit 
instanton candidates lie in $\ker(L_1)$, the space of harmonic 1-chains:
\begin{equation}
  \ker(L_1) = \ker(d_1) \cap \ker(\delta_1) \cong H_1(K; \R).
  \label{eq:kerL1}
\end{equation}

Moreover, if $A$ is the discrete audit 1-cochain, then the projection of 
$A$ onto $\ker(L_1)$ is non-zero precisely when there exists at least one 
audit instanton candidate in $K$:
\begin{equation}
  \|P_{\ker(L_1)} A\| > 0 \iff \exists \text{ at least one charged audit instanon candidate}.
  \label{eq:proj}
\end{equation}
\end{proposition}

\bigskip
\noindent
\textit{\textbf{Proposition \ref{prop:spectral}} $L_1$  $\cM$  $K$ 1-LaplacianInstanton1- $\ker(L_1)$1-furthermore$A$  $\ker(L_1)$ if and only if $K$ there existsInstanton}

% ===========================================================================
%  SECTION 9: DISCUSSION AND OPEN PROBLEMS
% ===========================================================================
\section{Discussion and Open Problems}
\section*{Open Problem}

\subsection{Summary / }

We have introduced \textbf{audit instantons} as a rigorous mathematical 
framework for detecting non-perturbative, topologically protected expert 
consensus failures in the SCX auditing paradigm. The central insight — 
that pointwise audit is structurally analogous to perturbative expansion, 
and thus systematically misses globally correlated topological defects — 
leads to a concrete detection protocol using persistent homology, with 
provable guarantees.

The key contributions are:
\begin{enumerate}[label=(\arabic*)]
  \item A cohomological definition of audit instantons (Definition \ref{def:instanton});
  \item A structural theorem linking homology classes to audit instanton 
        candidates (Theorem \ref{thm:main});
  \item A proof of invisibility to pointwise Yajie consensus (Theorem \ref{thm:yajie_blindness});
  \item An implementable detection protocol (Protocol \ref{prot:aid}).
\end{enumerate}

\bigskip
\noindent
\textit{we introduce\textbf{Instanton}SCXCore——thereforeMissing——Proof}

% ---------------------------------------------------------------------------
\subsection{Relation to the String Theory Analogy}
\subsection*{}

It is important to clarify the \emph{nature} of the analogy with string 
theory worldsheet instantons. We do \textbf{not} claim a formal 
mathematical correspondence between audit instantons and worldsheet 
instantons. Rather, the analogy operates at the level of \textbf{structural 
principles}:

\begin{center}
\begin{tabular}{@{}lll@{}}
\toprule
\textbf{Concept} & \textbf{String Theory} & \textbf{SCX Audit} \\
\textit{} & \textit{} & \textit{SCX} \\
\midrule
Non-perturbative effect & Worldsheet instanton & Audit instanton \\
\textit{} & \textit{Instanton} & \textit{Instanton} \\
\midrule
Perturbative expansion & Genus expansion & Pointwise Yajie audit \\
\textit{} & \textit{} & \textit{Yajie} \\
\midrule
Suppression factor & $\exp(-A/\alpha')$ & $\exp(-C/\sigma^2)$ \\
\textit{} & $\exp(-A/\alpha')$ & $\exp(-C/\sigma^2)$ \\
\midrule
Topological classification & $\pi_n$ or $H_n$ & $H_1$ (persistent) \\
\textit{} & $\pi_n$  $H_n$ & $H_1$ \\
\midrule
Instanton number & Winding / Chern class & Holonomy $\Phi(\gamma)$ \\
\textit{Instanton} & \textit{ / Chern} & \textit{ $\Phi(\gamma)$} \\
\midrule
Detection & Moduli space geometry & Persistent homology \\
\textit{} & \textit{} & \textit{} \\
\bottomrule
\end{tabular}
\end{center}

The analogy is \textbf{heuristic in origin but rigorous in execution}: it 
motivated the search for a topological characterization of non-local audit 
failures, and the resulting mathematical framework stands independently of 
string theory.

\bigskip
\noindent
\textit{Instanton\textbf{}\textbf{}InstantonInstantonthere exists\textbf{}\textbf{Heuristic}}

% ---------------------------------------------------------------------------
\subsection{Open Problems}
\subsection*{Open Problem}

We identify several open problems for future investigation:

\begin{enumerate}[label=\textbf{OP\arabic*}]
  \item \textbf{Sharpness of bounds / }: Can the lower bound 
        $\eta$ in Theorem \ref{thm:main} be sharpened using more refined 
        geometric data (curvature of $\cM$, higher-order statistics of 
        expert deviations)?
  
  \item \textbf{Higher-dimensional audit instantons / Instanton}: 
        Can the definition be extended to $k$-cycles for $k \geq 2$? A 
        ``2-audit instanton'' would correspond to a 2-cycle (void/bubble) 
        in the density filtration, representing a region where expert 
        consensus \emph{encloses} a systematically wrong prediction volume.
  
  \item \textbf{Dynamical audit instantons / Instanton}: In the 
        time-dependent SCX framework (with evolving Spring gatekeeper 
        and drifting experts), do audit instantons exhibit dynamics — 
        creation, annihilation, merging? What is the analog of the 
        ``instanton gas'' approximation?
  
  \item \textbf{Statistical significance / }: Develop 
        hypothesis tests for the null hypothesis ``the holonomy is 
        consistent with zero'' vs.\ the alternative ``there exists an 
        audit instanton,'' accounting for the adaptive nature of the 
        cycle selection.
  
  \item \textbf{Algorithmic efficiency / }: Can the detection 
        be accelerated by first computing the 1-Laplacian eigenmodes 
        and restricting persistence computation to the subspace of 
        near-harmonic cycles (the ``topological subspace'')?
  
  \item \textbf{Empirical validation / }: Apply Protocol 
        \ref{prot:aid} to real SCX audit datasets (AlN MLIP, PDF 
        reconstruction, ACE energy surfaces) and measure the prevalence 
        and impact of audit instantons compared to isolated pointwise 
        expert disagreements.
  
  \item \textbf{Connection to adversarial robustness / robustness}: 
        Are audit instantons related to adversarial examples in deep 
        learning? Both involve low-data-density regions and systematic 
        model failures. Can audit instanton detection be repurposed as 
        an adversarial robustness diagnostic?
  
  \item \textbf{Connection to Morse theory / Morse}: 
        If $\rho$ is treated as a Morse function on $\cM$, the birth and 
        death of 1-cycles in the persistence diagram correspond to 
        critical points of $\rho$. Can the ``instanton action'' 
        $\sim \exp(-1/\rho)$ be given a Morse-theoretic interpretation?
\end{enumerate}

\bigskip
\noindent
\textit{we point out thatOpen Problem(OP1) (OP2) Instantonk-$k \ge 2$(OP3) Instanton(OP4) (OP5) (OP6) SCX(OP7) robustness(OP8) Morse}

% ===========================================================================
\section{Conclusion}
\section*{Conclusion}

\begin{bilingual}{Conclusion}{Conclusion}
We have introduced the concept of \textbf{audit instantons} — non-perturbative, 
topologically protected failures of expert consensus — and provided a 
rigorous mathematical framework for their definition, detection, and 
interpretation within the SCX auditing paradigm. The framework builds on 
persistent homology of the data density filtration and discrete Hodge 
theory, with structural inspiration from worldsheet instantons in string 
theory.

Audit instantons reveal a fundamental limitation of pointwise consensus 
auditing: by treating each data point independently, it systematically 
misses globally correlated failure modes concentrated in low-density 
regions of the data manifold. These failures are \emph{topologically 
robust} — they cannot be eliminated by adding more auditors or adjusting 
consensus thresholds, but require filling the topological ``holes'' in 
the data distribution.

The detection protocol we provide is concrete, algorithmically 
implementable, and computationally tractable for realistic SCX 
deployments. We believe audit instantons open a new direction in 
trustworthy AI auditing — one that moves beyond local perturbation 
theory to a genuinely topological understanding of when and why expert 
systems fail.

\bigskip
\noindent
\textit{we introduce\textbf{Instanton}————SCXdefineHodgeInstanton}

\textit{InstantoneachMissing\textbf{}——``''}

\textit{SCXInstantonAI——Why}
\end{bilingual}

% ===========================================================================
%  ACKNOWLEDGMENTS
% ===========================================================================
\section*{Acknowledgments / Acknowledgments}

The authors thank the SCX research community at Xiaogan Supercomputing 
Center for insightful discussions on the topology of expert auditing. 
The structural analogy with string theory worldsheet instantons was 
developed during the ``String Theory Mathematics for SCX'' exploration 
\cite{scx_string_exploration}. We acknowledge the GUDHI and Ripser 
development teams for providing the persistent homology computational 
infrastructure on which the detection protocol depends.

\bigskip
\noindent
\textit{SCXXiaoganSCXInstanton``SCX''GUDHIRipser}

% ===========================================================================
%  APPENDIX: PROOF DETAILS
% ===========================================================================
\appendix
\section{Proof of the Holonomy Lower Bound}
\section*{lower boundProof}

We provide a more detailed proof of the isoperimetric inequality used in 
Theorem \ref{thm:main}, showing that $\oint_{\gamma} A > \eta$ under the 
positivity condition.

\begin{lemma}[Isoperimetric Inequality on $\cM$ / $\cM$ awaitingawaiting]
\label{lem:isoperimetric}
Let $\cM$ be a compact Riemannian $n$-manifold with Ricci curvature bounded 
below by $-(n-1)\kappa^2$ ($\kappa \geq 0$). For any 1-cycle $\gamma$ 
representing a non-trivial homology class $[\gamma] \in H_1(\cM; \Z)$, there 
exists a constant $C_{\cM} > 0$ depending only on the geometry of $\cM$ 
such that for any 2-chain $\Sigma$ with $\partial \Sigma = \gamma$:
\begin{equation}
  \vol_{n-1}(\Sigma) \geq C_{\cM} \cdot \frac{\left|\oint_{\gamma} A\right|}
    {\|A\|_{L^{\infty}} \cdot \length(\gamma)} \cdot \vol_n(\text{interior supported region}).
  \label{eq:isop}
\end{equation}
\end{lemma}

\begin{proof}
By the coarea formula and the fact that $A = d_0 \bar{f}$, we have:
\begin{equation}
  \oint_{\gamma} A = \int_{\gamma} d_0 \bar{f} = \int_{\Sigma} d_1 d_0 \bar{f} = 0,
  \label{eq:coarea_attempt}
\end{equation}
if $\Sigma \subset \cM$. This apparent contradiction is resolved by noting 
that \emph{no} 2-chain $\Sigma$ in $\cM_{\rho_{\text{crit}}}$ bounds 
$\gamma$ — indeed, condition (AI3) asserts precisely that $\gamma$ is not 
a boundary in the low-density region.

For any 2-chain $\Sigma \subset \cM$ with $\partial \Sigma = \gamma$, 
$\Sigma$ must intersect $\cM \setminus \cM_{\rho_{\text{crit}}}$ (the 
high-density region). Applying the isoperimetric inequality on $\cM$ 
\cite{chavel2001}, the volume of $\Sigma$ is bounded below by the 
\textbf{filling area} of $\gamma$, which for non-trivial homology 
classes scales with the injectivity radius of $\cM$ and the homological 
systole.

Explicitly, let $\text{sys}_1(\cM_{\rho_{\text{crit}}})$ denote the 
1-systole of the low-density region (the minimal length of a non-trivial 
1-cycle in $\cM_{\rho_{\text{crit}}}$). Then any bounding chain $\Sigma$ 
must have $(n-1)$-volume at least $\text{sys}_1(\cM_{\rho_{\text{crit}}})$, 
and the holonomy bound follows from the positivity condition 
(\ref{eq:positivity}), which gives a pointwise lower bound on the 
integrand of $\oint_{\gamma} A$.
\end{proof}

\bigskip
\noindent
\textit{\textbf{lemma \ref{lem:isoperimetric}$\cM$ awaitingawaiting} $\cM$  Ricci lower bound $-(n-1)\kappa^2$  Riemann $n$- $[\gamma] \in H_1(\cM; \Z)$ 1- $\gamma$there exists $\cM$  $C_{\cM} > 0$such thatsatisfying $\partial \Sigma = \gamma$ 2- $\Sigma$lower bound}

% ---------------------------------------------------------------------------
\section{Discrete Formulation for Numerical Implementation}
\section*{}

For readers implementing Protocol \ref{prot:aid}, we provide the discrete 
analogues of all key quantities.

% (We present the English version in the main text; the Chinese version 
%  appears as inline notes for conciseness.)

Let $K$ be a finite simplicial complex with:
\begin{itemize}
  \item Vertices $V = \{v_1, \ldots, v_N\}$ (audit data points);
  \item Edges $E = \{e_1, \ldots, e_{|E|}\}$ (pairs of nearby points);
  \item Triangles $T = \{t_1, \ldots, t_{|T|}\}$ (triples, for $n \geq 2$).
\end{itemize}

\bigskip
\noindent
\textit{ $K$  $V$ $E$  $T$}

\medskip
\noindent
\textbf{Discrete boundary operators / }:
\begin{align}
  (\partial_1)_{ve} &= \pm 1 \text{ if } v \in \partial e \text{ (with consistent orientation)}, \\
  (\partial_2)_{et} &= \pm 1 \text{ if } e \in \partial t.
\end{align}

\noindent
\textbf{Discrete 1-Laplacian / 1-Laplacian}:
\begin{equation}
  L_1 = \partial_2 \partial_2^T + \partial_1^T \partial_1.
  \label{eq:disc_laplacian}
\end{equation}

\noindent
\textbf{Discrete audit 1-cochain / 1-}:
\begin{equation}
  A \in \R^{|E|}, \quad A(e = [u, v]) = \bar{f}(v) - \bar{f}(u).
  \label{eq:disc_A}
\end{equation}

\noindent
\textbf{Discrete holonomy / }:
For a 1-cycle $\gamma = \sum_{e \in E} c_e e$ with integer coefficients 
$c_e$, the holonomy is:
\begin{equation}
  \Phi(\gamma) = \sum_{e \in E} c_e \cdot A(e) = c^T A.
  \label{eq:disc_hol}
\end{equation}

\noindent
\textbf{Harmonic projection detection / }:
Compute the eigenvectors $\{u_j\}_{j=1}^{\dim \ker(L_1)}$ of $L_1$ with 
eigenvalue zero. Then:
\begin{equation}
  \|P_{\ker(L_1)} A\|^2 = \sum_{j=1}^{\dim \ker(L_1)} |\langle u_j, A \rangle|^2.
  \label{eq:disc_proj}
\end{equation}
A non-zero norm signals the presence of at least one audit instanton candidate.

\bigskip
\noindent
\textit{ $L_1$  $\{u_j\}$ $\|P_{\ker(L_1)} A\|^2 = \sum_j |\langle u_j, A \rangle|^2$Instantonthere exists}

% ===========================================================================
%  BIBLIOGRAPHY
% ===========================================================================
\begin{thebibliography}{99}

\bibitem{scx_framework}
SCX Research Group.
\textit{SCX: Situs Consensus eXpert — A Discrete Hodge-Theoretic Framework for 
Multi-Expert Auditing.}
Xiaogan Supercomputing Center Technical Report, 2025--2026.

\bibitem{scx_hodge}
SCX Research Group.
\textit{Discrete Hodge Theory in SCX: Foundations and Applications.}
Xiaogan Supercomputing Center Technical Report, 2026.

\bibitem{scx_string_exploration}
SCX Research Group.
\textit{Can String Theory Mathematics Enrich SCX? — An Honest Creative Exploration.}
Xiaogan Supercomputing Center Internal Document, July 2026.

\bibitem{coleman1977}
S. Coleman.
\textit{The Uses of Instantons.}
In: \emph{The Whys of Subnuclear Physics} (Erice 1977), pp.\ 805--916.
Plenum Press, 1979.

\bibitem{polyakov1977}
A.\ M.\ Polyakov.
\textit{Quark Confinement and Topology of Gauge Groups.}
Nucl.\ Phys.\ B \textbf{120}, 429--458 (1977).

\bibitem{zomorodian2005}
A.\ Zomorodian and G.\ Carlsson.
\textit{Computing Persistent Homology.}
Discrete Comput.\ Geom.\ \textbf{33}, 249--274 (2005).

\bibitem{edelsbrunner2008}
H.\ Edelsbrunner and J.\ Harer.
\textit{Persistent Homology — a survey.}
In: \emph{Surveys on Discrete and Computational Geometry}, pp.\ 257--282.
AMS, 2008.

\bibitem{chazal2017}
F.\ Chazal and B.\ Michel.
\textit{An Introduction to Topological Data Analysis: Fundamental and Practical 
Aspects for Data Scientists.}
arXiv:1710.04019 (2017).

\bibitem{bott_tu}
R.\ Bott and L.\ Tu.
\textit{Differential Forms in Algebraic Topology.}
Graduate Texts in Mathematics, Vol.\ 82.
Springer, 1982.

\bibitem{bauer2021}
U.\ Bauer.
\textit{Ripser: Efficient Computation of Vietoris--Rips Persistence Barcodes.}
J.\ Appl.\ Comput.\ Topol.\ \textbf{5}, 391--423 (2021).

\bibitem{gudhi}
The GUDHI Project.
\textit{GUDHI User and Reference Manual.}
3.9.0 edition, 2023.
\url{https://gudhi.inria.fr/}

\bibitem{chavel2001}
I.\ Chavel.
\textit{Isoperimetric Inequalities: Differential Geometric and Analytic 
Perspectives.}
Cambridge University Press, 2001.

\bibitem{witten1982}
E.\ Witten.
\textit{Supersymmetry and Morse Theory.}
J.\ Diff.\ Geom.\ \textbf{17}, 661--692 (1982).

\bibitem{nakahara}
M.\ Nakahara.
\textit{Geometry, Topology and Physics.}
2nd edition. CRC Press, 2003.

\bibitem{hatcher}
A.\ Hatcher.
\textit{Algebraic Topology.}
Cambridge University Press, 2002.

\bibitem{green_string}
M.\ B.\ Green, J.\ H.\ Schwarz, and E.\ Witten.
\textit{Superstring Theory}, Vol.\ 1 \& 2.
Cambridge University Press, 1987.

\bibitem{polchinski}
J.\ Polchinski.
\textit{String Theory}, Vol.\ 1 \& 2.
Cambridge University Press, 1998.

\bibitem{douglas2006}
M.\ R.\ Douglas and S.\ Kachru.
\textit{Flux Compactification.}
Rev.\ Mod.\ Phys.\ \textbf{79}, 733--796 (2007).

\end{thebibliography}

% ===========================================================================
\end{document}
